\chapter{Conclusion}\label{chap:conclusion}

\section{Summary of Key Findings}
The implementation of counterfactual analysis on the path-based knowledge graph
recommender system has yielded significant insights into the explainability of
AI-driven recommendations. Our research demonstrates that by examining alternative
scenarios and attribute modifications, we can shed light on the decision-making
process of complex recommender systems. This approach allows us to move beyond the
"black box" nature of many AI systems and provide clearer explanations for why specific
products are recommended to users.

A key finding of our study is the identification of influential attributes and path
types in the recommendation process. By analyzing different path structures, such
as \texttt{'mentions' -> 'rev\_described\_by'} and \texttt{'purchase' -> 'also\_viewed'
-> 'rev\_also\_bought'}, we have uncovered the underlying logic driving the
recommender system's choices. This granular understanding of path influences enhances
the system's transparency, allowing both developers and users to better
comprehend the factors contributing to each recommendation.

Moreover, our counterfactual analysis revealed the relative importance of various
product attributes, including brand associations, descriptive keywords, and
product relationships. By systematically altering these attributes and observing
the resulting changes in recommendation scores, we have developed a framework
for explaining why certain products are recommended over others. This approach not
only improves the explainability of the recommender system but also provides valuable
insights into user preferences and decision-making patterns. However, it is important
to note the limitations of our method, such as the computational complexity and
the assumption of attribute independence, which present opportunities for future
research to further refine and expand upon this explainable AI approach in recommender
systems.

\section{Contributions to Explainable AI in Recommender Systems}
Our methodology makes several key contributions to the field of explainable AI in
recommender systems. First, it provides a novel approach to generating "what-if"
scenarios, allowing for a more comprehensive understanding of the system's
decision boundaries. This goes beyond traditional explanation methods that focus
solely on the attributes present in the recommended item.

Second, our approach offers a more nuanced view of feature importance in
recommendations. By identifying which attributes could lead to the same recommendation,
we provide insights into the relative influence of different features, enhancing
the interpretability of the system's decisions.

\section{Marketing and Business Implications}
From a marketing perspective, our analysis reveals valuable insights for product
positioning and targeted advertising. The diverse range of brands and categories
associated with recommended products suggests opportunities for flexible marketing
strategies. Furthermore, the identification of key descriptive words and related
products opens avenues for content marketing and cross-promotion strategies.

\section{Future Research Directions}
Future research in the domain of path-based knowledge graph recommender systems could
explore the integration of scenarios that involve multiple counterfactual attributes
simultaneously. This approach would address the complexity of real-world
decision-making processes where multiple factors are often intertwined, affecting
user preferences and product recommendations concurrently. By analyzing the
combined effects of various attributes, researchers could gain a more nuanced understanding
of their interdependencies and the synergistic effects on recommendation outputs.
Additionally, dynamic alterations of embeddings represent a promising area for enhancing
the adaptiveness and accuracy of counterfactual explainability of recommender systems.
Such efforts would advance the explainability of these systems, enabling clearer
insights into how and why recommendations change under varied conditions.